% Technical appendix (companion) for the Chain-NPM PVLDB submission.
%
% Purpose: hold the material the 12-page submission cannot carry, so that the
% submission shrinks without losing evidence.  Every item here is cited from
% the corresponding place in the submission; the moved blocks are the
% submission's own text, with cross-references resolved through xr.
%
% This file lives next to main.tex so that it uses the VLDB-bundled
% acmart.cls and the same refs.bib / figs/.
%
% Build order: main.tex first (its main.aux is imported by xr), then
%   latexmk -xelatex tech-appendix.tex

\documentclass[sigconf, nonacm]{acmart}

\usepackage{booktabs}
\usepackage{amsmath}
\usepackage{xspace}
\usepackage{algorithm}
\usepackage{algpseudocode}
\usepackage{xr}
\externaldocument[P-]{main}

\setlength{\emergencystretch}{2em}

\newcommand{\chainnpm}{\textsc{Chain-NPM}\xspace}
\newcommand{\privlava}{\textsc{PrivLava}\xspace}
\newcommand{\privpetal}{\textsc{PrivPetal}\xspace}
\newcommand{\privmrf}{\textsc{PrivMRF}\xspace}
\newcommand{\privmrfden}{\privmrf(denorm)\xspace}

% Cross-document references: P-<label> is the submission's label imported by
% xr, so \pref{tab:tpch} prints the submission's number for that table.
\newcommand{\pref}[1]{\ref{P-#1}}
\newcommand{\psec}[1]{Section~\ref{P-sec:#1} of the submission}
\newcommand{\pfig}[1]{Figure~\ref{P-fig:#1} of the submission}
\newcommand{\ptab}[1]{Table~\ref{P-tab:#1} of the submission}
\newcommand{\pprop}[1]{Proposition~\ref{P-prop:#1} of the submission}

\title{Chain-NPM: Technical Appendix}
\subtitle{Supplementary material for the submitted paper}

\author{Sheng Xiang}
\affiliation{%
  \institution{Tongji University}
}
\email{xiangsheng@tongji.edu.cn}

\author{Yunpeng Wang}
\affiliation{%
  \institution{Tongji University}
}
\email{2251739@tongji.edu.cn}

\author{Weihao Wang}
\affiliation{%
  \institution{Tongji University}
}
\email{2634137@tongji.edu.cn}

\author{Ying Zhang}
\affiliation{%
  \institution{Zhejiang Gongshang University}
}
\email{ying.zhang@zjgsu.edu.cn}

\begin{document}

\begin{abstract}
This companion holds the material that the page-limited submission cannot
carry: the witness table for the identifiability barrier, the Phase-2
walk-through, the numeric tables behind the controlled-benchmark figure, the
deep TPC-H and ACS rails, the planted-effect and natural-mode tables, the
three real-benchmark rails figure, and the budget-allocation, ancestor-scale
and cost--effect ablations.  The submission cites each item where it is
relevant; nothing here adds a new claim.  Section, figure, table and
proposition numbers written as ``of the submission'' refer to the paper and
are resolved at build time.
\end{abstract}

\maketitle

%=======================================================================
\section{Sensitivity constants and total-variation statements}
\label{app:proofs}

\begin{proof}[Proof of Lemma~\pref{lem:sens}]
Fix neighboring databases $D,D'$ differing in a single individual tuple
$i^\star$ of group $g^\star$ (size $s=\min(|g^\star|,\tau)\ge1$). Every
released marginal is a sum over groups of a per-group normalized
contribution, and only the $g^\star$ term differs.
\emph{(i) Intra-table 1-way.} The affected count changes by $1$, and the
public-free normalization denominator by $1$; the $\ell_2$ change of the
vector is $\le1$.
\emph{(ii) Adjacent- and cross-hop group-normalized marginals.} The
per-group term is $c_g=\frac{1}{s_g}\sum_{i\in g}x_i$ with $x_i$ a one-hot
attribute vector. Adding $i^\star$ (with $s<\tau$) changes $c_{g^\star}$ by
$\|x-c\|_1/(s+1)\le2/(s+1)\le1$; with $s=\tau$ the clipped count changes by
$\le1/\tau\le1$; removing $i^\star$ changes it by $\|x-c\|_1/s\le2/s$, which
equals $1$ when $s=1$ (the group vanishes) and is $\le1$ otherwise. Hence
the $\ell_2$ sensitivity is $\le1$.
\emph{(iii) Within-group pair marginals.} The per-group pair term changes
only through pairs involving $i^\star$: at most $2(s-1)/s^2\le1$ in mass,
plus a renormalization term $\le1$ when $s$ crosses the truncation; the
worst case ($s=2$ removal) gives $\ell_2$ change $\le2$.
Since only one group term moves in each case, the stated bounds are the
$\ell_2$ sensitivities.
\end{proof}

\begin{lemma}[Tight $\ell_2$ sensitivities]
\label{lem:sens-tight}
Under the one-individual adjacency of Lemma~\pref{lem:sens}, let $s$ be the
truncated size of the affected group. Then:
(i) intra-table 1-way marginals have $\Delta_2=1$;
(ii) adjacent- and cross-hop group-normalized marginals have
$\Delta_2=1$, attained by removal of a single-individual group, while any
addition satisfies $\Delta_2\le2/(s+1)<1$;
(iii) within-group pair marginals have $\Delta_2=2$ in the statement of
Lemma~\pref{lem:sens}, with the maximizing case a removal at $s=2$; for
$s\ge3$ the per-group change is at most $2(s-1)/s^2\le3/4$ plus a
$O(1/s^2)$ renormalization term.
\end{lemma}
\begin{proof}
The three cases are precisely the per-group computations of
Lemma~\pref{lem:sens}; the values above record the supremum over $s$ and
over add/remove together with explicit witnesses (single-individual group
removal for (i)--(ii), two-individual group removal for (iii)), hence the
bounds are tight.
\end{proof}

\begin{corollary}[Total variation of released marginals]
\label{cor:tv}
Let $M\in\mathbb{R}^d$ be any marginal of Lemma~\pref{lem:sens} with
$\ell_2$ sensitivity $\Delta_2$, released as
$\hat M=M+\mathcal{N}(0,\sigma^2 I)$. With probability $\ge1-\beta$,
$\|\hat M-M\|_\infty\le\sigma\sqrt{2\ln(2d/\beta)}$; if $M$ is normalized
by a group count of magnitude $n_1 m_{\min}$, then
$\mathrm{TV}(\hat M/\|\hat M\|_1,\;M/\|M\|_1)
\le \frac{d\,\sigma}{n_1 m_{\min}}\sqrt{2\ln(2d/\beta)}$.
\end{corollary}
\begin{proof}
A union bound over the $d$ Gaussian coordinates gives the
$\ell_\infty$ statement. Total variation is at most half the $\ell_1$
difference of the normalized vectors, hence at most
$\frac{d}{2}\|\cdot\|_\infty$; dividing by a mass of magnitude
$n_1m_{\min}$ converts absolute to relative error, absorbing the
normalization perturbation into the same order.
\end{proof}

\begin{lemma}[TV decay along the chain of posteriors]
\label{lem:tvdecay}
Consider a $k$-hop chain with true channels $P_{Z_i|Z_{i-1}}$ and
synthesized channels $\hat P_{Z_i|Z_{i-1}}$ satisfying
$\max_{z_{i-1}}\mathrm{TV}(\hat P_{Z_i|z_{i-1}},\,P_{Z_i|z_{i-1}})
\le\eta_i$. Then the synthesized endpoint joint satisfies
$\mathrm{TV}(\hat P_{Z_0,Z_k},\,P_{Z_0,Z_k})
\le 1-\prod_{i=1}^k(1-\eta_i)\le\sum_{i=1}^k\eta_i$, and the recoverable
endpoint correlation obeys
$|\Delta_{\mathrm{syn}}|\le\prod_i(1-\eta_i)\,\Delta_{\mathrm{real}}
+2\sum_i\eta_i$. Each $\eta_i$ is lower-bounded by the minimax TV risk of
estimating the hop's channel from its budget slice under the Gaussian
mechanism, hence is bounded away from $0$ at any fixed total budget.
\end{lemma}
\begin{proof}
Write the joint as a Markov chain and telescope: replacing the $i$-th
channel changes the joint distribution by at most $\eta_i$ times the mass
that reaches hop $i$, which is at most the product of the survival
factors $(1-\eta_j)$ of the preceding hops plus the channel error itself;
induction yields the product form, and the data-processing inequality for
TV gives the additive relaxation. The correlation bound follows by
writing $\Delta$ as a linear functional of the endpoint joint and
bounding by TV. The minimax claim is the standard Gaussian-mean lower
bound (Le Cam's method applied to the channel entries at the hop's
$\gamma_i$ slice).
\end{proof}

%=======================================================================
\section{Planting protocol, V-shaped sensitivity, and statistical conventions}
\label{app:plant}

\textbf{Planting protocol.} To instantiate the barrier of
Proposition~\pref{prop:barrier} on real schemas, we plant cross-hop
structure while leaving every one-hop object untouched. On a chain
$R_0{\leftarrow}R_1{\leftarrow}R_2$ we pick two balanced root anchors
$X_1,X_2$ and set the latent $Z=X_1\oplus X_2$ (pure interaction, no main
effects) under \emph{plant\_xor}, or $Z=X_1$ under \emph{plant\_product}.
The leaf response $A$ is rewritten
with probability $p$ by the rule ($A\sim\mathrm{Bern}(r_Z)$ or
$A=U\oplus Z$) and kept real otherwise; then, within each value class of
a middle attribute, count-preserving swaps restore the
$(R_1\text{-attr}\times A)$ contingency exactly. Hence the FK
cardinalities and every one-hop marginal the per-FK family consumes are
unchanged by construction, while the cross-hop functional moves by a
calibrated amount; $p$ is chosen by bisection over seeds so that the
planted effect hits $\Delta\in\{0.10,0.20,0.316,0.40\}$ under
\emph{plant\_product}, or the XOR target value $0.8$ under
\emph{plant\_xor}. Repeating the protocol with the rule family
$\{\oplus,\ \bmod k,\ \mathbf 1[\cdot\ge t],\ \vee\}$ forms the
non-linearity scan.

\textbf{V-shaped sensitivity.} For a star $R_1{\leftarrow}C{\to}R_2$, the
star workload adds the cross-parent marginal $R_1{\times}R_2$ and the
triple $R_1{\times}R_2{\times}C$, each child row weighted by the inverse
of its $R_1$-group size (truncated at $\tau$). For a neighbor obtained by
adding or removing one $R_1$ tuple together with its $C$ rows, each star
query's $\ell_2$ sensitivity is at most $1$ by the per-group argument of
Lemma~\pref{lem:sens}; the $R_2$-side neighbor follows by the same
argument with roles swapped.

\textbf{Statistical conventions.} Ten seeds per configuration with common
random numbers (data, query-workload, and planting seeds shared across
methods); medians with bootstrap 95\% confidence intervals; per-seed
paired Wilcoxon signed-rank tests with Holm correction over the family
``\chainnpm{} vs.\ each baseline'' per dataset, $\varepsilon$, and metric.
Seeds,
$\varepsilon$ grids, and metric definitions are fixed before baseline
tuning and released with the artifact.

%=======================================================================
\section{The deep-chain normalization study}
\label{app:normalization}

Measuring a cross-hop marginal is not sufficient on its own; two normalization
choices matter, and both show up as failures of effect retention. (i) The
cross-hop marginal must be normalized \emph{once}, by the FK group it is
traversed through. Normalizing a second time by the ancestor group collapses
the released cross-hop mass from $54{,}369$ to $1{,}000$ while the noise is
unchanged, leaving the measured object at signal-to-noise ratio $O(1)$---and
the resulting retention error returns to $1.145$, i.e.\ \emph{worse} than not
measuring the marginal at all ($1.027$), against $0.380$ under the rule the
paper adopts. (ii) The lift denominator is itself a noised one-way marginal of
a small table, so an unguarded cross-hop lift manufactures a spurious contrast
of $+0.15$--$+0.20$ at $\varepsilon\le0.4$ where the true contrast is
$\approx-0.004$; the existing SNR shrinkage cuts that artefact by a factor of
$1.9$--$4.9$ at $\varepsilon{=}0.2$--$0.8$. It is not a cure everywhere: at
$\varepsilon{=}0.1$ the guard barely moves the artefact
($+0.197\to+0.177$), because the noised root marginal sits far above the noise
scale, and the natural-mode two-hop error there also worsens
($3.88\to4.34$, both measured without the quota calibration), so at
$\varepsilon{=}0.1$ the cross-hop object costs more than it returns. The
row-level quota calibration of \pref{sec:setup} addresses exactly this cell:
the natural two-hop error drops from $4.34$ to $0.28$ without changing a
released marginal. With both rules the planted effect on the deep chain is
recovered ($1.027\to0.380$), and the guard is inert on the other three
datasets, whose V-shaped path already measured its cross-parent object. Two
further costs of the rule are worth stating plainly: the twelve extra
cross-hop marginals raise the released-marginal count from $31$ to $43$ and,
under an equal budget split, the noise on every query by $17\%$; and where the
guard does not bite the extra marginals dilute accuracy, as the deep TPC-H
rail below shows.

\section{The identifiability barrier on a witness construction}
\label{app:barrier}

\pprop{barrier} is proved by a two-database witness that a reader can verify
row by row; the submission states the barrier and defers the table to this
companion.

\begin{table}[t]
\centering
\caption{A minimal instantiation of the identifiability barrier
(\pref{prop:barrier}). Two databases $\mathcal{D}^+$ and
$\mathcal{D}^-$ over the same three-level chain ($2$ regions \penalty0$\times$\penalty0
$2$ households \penalty0$\times$\penalty0 $2$ individuals per household, uniform group
sizes) share \emph{every} per-FK view---all table marginals are
uniform, $Z\perp U$ makes the $R_0{\to}R_1$ view coincide, and the
within-group composition of the $R_1{\to}R_2$ view is identical because
the sign of $A$ is symmetric in $U$ (the $Z{=}0$ households swap roles
between the two databases). Yet the two-hop join functional
$\varphi=\Pr[A=U\oplus Z]$ equals $1$ on $\mathcal{D}^+$ and $0$ on
$\mathcal{D}^-$. Any synthesizer whose input is the multiset of per-FK
views receives identical input in the two cases and therefore errs by at
least $1/2$ on $\varphi$ for one of them, at any budget.}
\label{tab:barrier}
\begin{tabular}{cccc}
\toprule
$U$ ($R_0$) & $Z$ ($R_1$) & $A$ in $\mathcal{D}^+$ & $A$ in $\mathcal{D}^-$ \\
\midrule
0 & 0 & 0 & 1 \\
0 & 0 & 0 & 1 \\
\cmidrule(lr){1-4}
0 & 1 & 1 & 0 \\
0 & 1 & 1 & 0 \\
\cmidrule(lr){1-4}
1 & 0 & 1 & 0 \\
1 & 0 & 1 & 0 \\
\cmidrule(lr){1-4}
1 & 1 & 0 & 1 \\
1 & 1 & 0 & 1 \\
\bottomrule
\end{tabular}
\end{table}

%=======================================================================
\section{Phase-2 walk-through}
\label{app:emwalk}

The submission describes the EM soft assignment and the ancestral sampler in
prose; the figure below traces one refinement-and-sampling step on the
running example's micro-instance.

\begin{figure}[t]
\centering
\includegraphics[width=\columnwidth]{figs/fig4_walk.pdf}
\caption{Walk-through of Phase~2 on the running example's micro-instance:
the noisy cross-hop marginal $\tilde f$ (left, Line~3) is fed to the EM
soft assignment of each cell $c$ to the latent household type
$z\in\{1,0\}$ with weights $w_{c,z}$ (middle), and the fitted conditionals
then drive ancestral synthesis $\hat r\to\hat h\to\hat\imath$ (right).
Both stages operate solely on the released noisy statistics---pure
post-processing---and therefore consume zero additional privacy budget
(Algorithms~\pref{alg:em}--\pref{alg:ancestral}).}
\label{fig:emwalk}
\end{figure}

%=======================================================================
\section{Numeric tables for the controlled multi-hop benchmark}
\label{app:controlled}

The submission reports the controlled benchmark of \psec{gap-evidence} as a
figure (\pfig{controlled}) and quotes the per-cell numbers in its prose.
The two tables below are the same runs and medians in tabular form.

\begin{table}[t]
\centering
\caption{Two-hop $\Delta$-preservation ratio $\Delta_{\mathrm{syn}}/\Delta_{\mathrm{real}}$.
Per-FK family stays within the $\pm 0.2$ per-run band of the no-cross-table
level at every $\varepsilon$; the joint reference is consistently positive.
The \privlava{} row is its released code; a faithful implementation of its
Algorithm~3 stays within $0.04$ of it in these cells.}
\label{tab:gap}
\begin{tabular}{lccc}
\toprule
Method & $\varepsilon{=}0.8$ & $\varepsilon{=}3.2$ & $\varepsilon{=}8.0$ \\
\midrule
\privmrf{} (lower bound) & $0.035$ & $0.002$ & $0.009$ \\
\privlava & $0.009$ & $0.038$ & $-0.011$ \\
\shortstack[l]{\privmrfden{}\\ (joint ref.)} & $\mathbf{0.205}$ & $\mathbf{0.239}$ & $\mathbf{0.248}$ \\
\bottomrule
\end{tabular}
\end{table}

\begin{table}[t]
\centering
\caption{Query relative error (large-count queries, $\varepsilon{=}3.2$),
reported per hop class ($1$-hop R$\Join$H, $1$-hop H$\Join$I, $2$-hop
R$\Join$H$\Join$I), for each method, with the two \chainnpm{} variants grouped
below the baselines. \privlava{} is its released code; a faithful
implementation of its Algorithm~3 stays within $0.12$ of it here.}
\label{tab:re32}
\begin{tabular}{lccc}
\toprule
Method & 1-hop RE & 1-hop RE & 2-hop RE \\
 & (R$\Join$H) & (H$\Join$I) & (R$\Join$H$\Join$I) \\
\midrule
\privmrf & $0.240$ & $0.117$ & $0.212$ \\
\shortstack[l]{\privmrfden{}\\ (joint ref.)} & $\mathbf{0.058}$ & $0.323$ & $0.404$ \\
\privlava & $0.285$ & $0.256$ & $0.190$ \\
\midrule
\chainnpm-direct & $0.105$ & $\mathbf{0.060}$ & $\mathbf{0.089}$ \\
\chainnpm-mixture & $0.098$ & $0.063$ & $0.095$ \\
\bottomrule
\end{tabular}
\end{table}

%=======================================================================
\section{TPC-H rails (standard and deep)}
\label{app:tpch-rails}

\textbf{Standard rail: what TPC-H does and does not separate.} \label{app:tpch}

We use the TPC-H benchmark (SF=1) with the correlation-planting
transformations of~\cite{cai2023privlava} (adopted by
PrivPetal~\cite{cai2025privpetal}): (i) lineitem prices/quantities are
scaled by part-type and part-brand weights; (ii) each year's orders are
concentrated on a random 15\% subset of customers, coupling
\emph{customer nation} to \emph{order year}; (iii) orders are dropped
unless their size matches the order year, coupling \emph{order size} to
\emph{order year}. Private relations: Orders and Lineitem; public:
Customer, Nation, Part. This yields 536K orders and 2.1M lineitems.
Ground truth for evaluation is the planted database converted back to the
original schema (identical to PrivPetal's own postprocessing).

\chainnpm runs on this schema with the nation attribute joined into Orders
(public joins are free), $Q{=}42$ joint chain marginals in one budget pass,
$\delta=1/n_{\mathrm{leaf}}$. PrivPetal is run with its official code
(tuple\_num$=$2, process\_num$=$4). Metrics: median relative error of
multi-hop conjunctive count queries (2-hop Nation$\Join$Orders$\Join$Lineitem
and 1-hop Orders$\Join$Lineitem), and preservation ratios of the planted
associations.

\noindent Additional planted-association preservation (ratios,
$\varepsilon{=}0.2$ / $\varepsilon{=}3.2$): order-size\penalty0$\leftrightarrow$\penalty0year
$1.02$ / $0.98$ (PrivPetal), $1.16$ / $1.16$ (\chainnpm); mean lineitem
quantity $1.01$ / $1.01$ (PrivPetal), $0.98$ / $0.96$ (\chainnpm).
PrivPetal shows noticeable run-to-run variance at low budget (a repeat
$\varepsilon{=}0.2$ run gave 2-hop RE $0.072$ and 1-hop RE $0.019$);
we report the clean rerun and note the stochasticity.

\paragraph{Discussion: what TPC-H does and does not separate.}
On the standard TPC-H setup, PrivPetal performs well---and so does
\chainnpm. This is expected and informative: with only \emph{two} private
relations (Orders, Lineitem) and all upstream tables public, every planted
correlation lives within a single FK pair after the public join (the
nation\penalty0$\leftrightarrow$\penalty0year association is carried by Orders' foreign key to
the public Customer table), so per-FK models \emph{can} express it. The two
methods are indeed comparable here (per-cell numbers in the
technical appendix): \chainnpm has a
slight edge on two-hop
query error at low budget ($0.044$ vs $0.083$), PrivPetal a slight edge at
high budget ($0.015$ vs $0.031$), and both preserve all planted
associations. The structural gap of Section~\ref{P-sec:limit-mech} instead requires
\emph{deep private chains}---as in the controlled three-table benchmark of
Section~\ref{P-sec:gap-evidence}, where the per-FK family's $\Delta$-recovery
stays at $\approx 0$ regardless of budget. The next paragraph tests this on
the real schema.

\paragraph{Deep-chain variant: \textsc{Customer} private.}
A deep variant with a \emph{three}-level private chain
$\textsc{Customer}\to\textsc{Orders}\to\textsc{Lineitem}$---the real-schema
instance of the structure Section~\ref{P-sec:limit-mech} identifies---is
reported in full in the technical appendix. \chainnpm answers $3$-hop
queries at $5.4\%$ median relative error at $\varepsilon{=}3.2$ against
PrivPetal's $4.0\%$, and pays for the cross-hop object at the low budget
($33.6\%$ against $12.6\%$): its twelve extra marginals dilute the same
single budget. Where it differs from the latent-chain baseline is cost and
coverage---$\approx 22$--$24$ seconds of synthesis on one CPU against
$1.7$--$2.3$ hours per budget on an A800 GPU, one budget pass over $46$
marginals rather than a three-way split, and all $150$K customers
synthesized, including the $50$K without orders, where PrivPetal emits only
the $99{,}847$ with orders.

\paragraph{Efficiency.}
\chainnpm's entire pipeline (one budget pass over $42$ joint marginals plus
post-processing synthesis of $536$K orders and $2.1$M lineitems) completes in
$\approx 2.2$ minutes on a single CPU. PrivPetal's official implementation
takes $1.8$--$2.1$ hours per budget on two A800 GPUs for the same schema---a
$\approx 50\times$ wall-clock gap, which the hardware asymmetry makes
conservative for \chainnpm{} while our exclusion of the evaluation from the
timed window shades it the other way---mirroring the single-pass-vs-iterative
contrast on the controlled benchmark and following directly from EM being
post-processing rather than a budget-consuming loop. The deep-chain variant
shows the same pattern: $1.7$--$2.3$ hours per budget on one A800 for
PrivPetal versus $\approx 22$--$24$ seconds for \chainnpm.
\label{app:tpch-deep}

\textbf{Standard rail (Orders and Lineitem private).}
\begin{table}[t]
\centering
\caption{TPC-H (planted, SF=1). $\Delta$-ratio $=$ preserved/true
planted association (1.0 $=$ perfect); RE $=$ median relative error of
multi-hop count queries (2-hop
Nation\penalty0$\times$\penalty0Orders\penalty0$\times$\penalty0Lineitem).}
\label{tab:tpch}
\small
\begin{tabular}{@{}lcccc@{}}
\toprule
 & \multicolumn{2}{c}{$\varepsilon{=}0.2$} & \multicolumn{2}{c}{$\varepsilon{=}3.2$} \\
\cmidrule(lr){2-3} \cmidrule(lr){4-5}
Method & 2-hop RE & nat\penalty0$\leftrightarrow$\penalty0yr & 2-hop RE & nat\penalty0$\leftrightarrow$\penalty0yr \\
\midrule
PrivPetal (official) & $0.083$ & $1.05$ & $0.015$ & $0.99$ \\
\chainnpm (ours) & $0.044$ & $1.04$ & $0.031$ & $1.01$ \\
\bottomrule
\end{tabular}
\end{table}

\textbf{Deep rail (\textsc{Customer} private).} We build a deep TPC-H variant with a \emph{three}-level private chain
$\textsc{Customer}\to\textsc{Orders}\to\textsc{Lineitem}$: Customer carries
nation, market segment, and account-balance-bin attributes; the existing
plantings (nation\penalty0$\leftrightarrow$\penalty0year across Customer\penalty0$\to$\penalty0Orders,
size\penalty0$\leftrightarrow$\penalty0year and part-ratio-scaled quantities across
Orders\penalty0$\to$\penalty0Lineitem) are retained, and a new within-Customer
segment\penalty0$\leftrightarrow$\penalty0nation association is planted, so the composition
segment\penalty0$\to$\penalty0nation\penalty0$\to$\penalty0year\penalty0$\to$\penalty0quantity spans all three private tables.
Evaluation adds 3-hop conjunctive counts
(Nation\penalty0$\times$\penalty0Segment\penalty0$\times$\penalty0Year\penalty0$\times$\penalty0Quantity-bin) over the 3-table
join. \chainnpm extends to a three-level chain ($Q{=}46$ marginals, one
budget pass, intermediate rows sampled once and conditioning both hops, by
the construction of \psec{beyond});
PrivPetal runs its official multi-stage pipeline (budget split
$0.5/0.25/0.25$ over its Customer--Orders, Orders--Lineitem, and
Partsupp--Lineitem stages).

\begin{table}[t]
\centering
\caption{Deep TPC-H (Customer private, SF=1). 3-hop RE $=$ median relative
error of Nation\penalty0$\times$\penalty0Segment\penalty0$\times$\penalty0Year\penalty0$\times$\penalty0Qbin join counts;
seg\penalty0$\leftrightarrow$\penalty0year is the three-table composition signal
(preserved/true, 1.0 $=$ perfect).}
\label{tab:tpch-deep}
\small
\begin{tabular}{@{}lcccc@{}}
\toprule
 & \multicolumn{2}{c}{$\varepsilon{=}0.2$} & \multicolumn{2}{c}{$\varepsilon{=}3.2$} \\
\cmidrule(lr){2-3} \cmidrule(lr){4-5}
Method & 3-hop RE & seg\penalty0$\leftrightarrow$\penalty0yr & 3-hop RE & seg\penalty0$\leftrightarrow$\penalty0yr \\
\midrule
PrivPetal (official) & $0.126$ & $1.12$ & $0.040$ & $0.98$ \\
\chainnpm (ours) & $0.336$ & $1.72$ & $0.054$ & $1.73$ \\
\bottomrule
\end{tabular}
\end{table}

Both methods handle the deep chain. PrivPetal's multi-stage latent
pipeline answers 3-hop queries at $4.0\%$ median relative error at
$\varepsilon{=}3.2$ ($0.126$ at $\varepsilon{=}0.2$) and preserves all
planted associations (ratios $0.98$--$1.12$). \chainnpm attains $0.054$ at
$\varepsilon{=}3.2$ but pays for the cross-hop object at the low budget
($0.336$): its twelve extra marginals dilute the same single budget
(\pref{sec:four-datasets}), so at $\varepsilon{=}0.2$ the latent-chain
baseline is ahead on query error while \chainnpm remains ahead on cost and
coverage. This contrast is honest and informative:
PrivPetal is itself a chain-aware model---its child stage is conditioned
on the synthetic parent rows, so composed correlations propagate---and the
structural collapse of \psec{limit-mech} is the failure mode of the
PrivLava-style per-FK family (Table~\ref{tab:gap}), not of
PrivPetal-style
latent chains. Where \chainnpm differs is cost and coverage: the deep rail's
synthesis completes in $\approx 22$--$24$ seconds on a single CPU versus
$1.7$--$2.3$ hours per budget for \privpetal{}'s own pipeline on an A800 GPU
($\approx 300\times$ on the bases stated in \ref{tab:perbasis-cost}); it
consumes one
budget pass over $46$ marginals rather than a three-way split; and it
synthesizes all $150$K customers, including the $50$K without orders,
whereas PrivPetal emits only the $99{,}847$ customers with orders
(customer-only segment counts then incur RE $0.33$).
\chainnpm's segment\penalty0$\leftrightarrow$\penalty0year ratio slightly exceeds $1$
because the absolute composition signal is small (total-variation $0.33$
across $35$ cells) and both noise and lift-composition inflate small
values; the 3-hop query errors are the primary evidence. Deep chains are
intrinsically harder at low budget: at $\varepsilon{=}0.2$ the single pass
must spread $\sigma \approx 194$ Gaussian noise over $46$ marginals whose
effective mass is set by the $100$K customer level rather than the $536$K
order level. Two of the composition ratios thereby leave the band we would
like---the mean-lineitem-quantity ratio reaches $1.46$ at
$\varepsilon{=}0.2$, against $1.00$ at $\varepsilon{=}3.2$---and the
synthesized order and lineitem counts shrink by $6$--$8\%$; the correlation
structure itself is retained. This is the same correlations-cheap,
counts-expensive pattern as on the controlled benchmark, seen from the other
side.

%=======================================================================
\section{ACS PUMS California}
\label{app:acs-rail}

\textbf{The rail.}

To verify that \chainnpm works beyond engineered benchmarks, we run it on
the ACS PUMS 2022 California microdata: $151$K occupied households and
$383$K persons (vacant units excluded), with $7$ discretized attributes on
each table (household type\slash tenure\slash income\slash structure, person
age\slash sex\slash education\slash employment\slash relationship\slash hours). This is a natural
two-level chain with heavy-tailed group sizes (max $20$) and strong
size--attribute coupling---the stress case for the size model.
\chainnpm uses $Q{=}74$ marginals ($7$ one-way per table, $5$ parent pairs,
$49$ parent$\times$child cells, $6$ size histograms) with the size-itself
attribute excluded from size conditioning; PrivPetal runs its official
Census pipeline.

\chainnpm answers 2-hop conjunctive counts at $\approx 20\%$ median
relative error (large counts) and person-level marginals at around $10\%$
at both budgets (the technical appendix tabulates the rail), with
household-size totals within
$7\%$ of truth; 0-hop
household marginals stay below $2\%$. This trails the latent baseline:
PrivPetal reaches $5.8$--$9.0\%$ 2-hop and $1.5$--$2.0\%$ 1-hop relative
error with nearly exact counts ($-0.7\%$), at the price of hours of GPU
time per budget. The score is informative. The controlled benchmark of
Section~\ref{P-sec:gap-evidence} isolates the \emph{structural} per-FK
failure; ACS has no planted correlations to lose, so both methods are
learning the same $7{\times}7$-attribute joint, and PrivPetal's
per-stage latent models simply fit a two-level join more efficiently than
our single-pass pairwise marginals---a budget-allocation effect, not a
structural one. Two-hop joins over a $7{\times}7$ attribute product with
only $151$K parent rows are noise-limited at these budgets, and the
contrast confirms the mechanism's limitations on skewed, heavy-tailed
microdata while an exactly-calibrated size model and a single-pass
pipeline retain their cost and simplicity advantages.
\chainnpm uses $Q{=}74$ marginals ($7$ one-way per table, $5$ parent pairs,
$49$ parent$\times$child cells, $6$ size histograms) with the size-itself
attribute excluded from size conditioning; PrivPetal runs its official
Census pipeline.

\begin{table}[t]
\centering
\caption{ACS PUMS California 2022. RE (\pref{sec:setup}) on random
conjunctive counts with true answer $\ge 50$ (2-hop: household\penalty0$\times$\penalty0person
over the join; 1-hop: person-only).}
\label{tab:acs}
\small
\begin{tabular}{@{}lcccc@{}}
\toprule
 & \multicolumn{2}{c}{$\varepsilon{=}0.8$} & \multicolumn{2}{c}{$\varepsilon{=}3.2$} \\
\cmidrule(lr){2-3} \cmidrule(lr){4-5}
Method & 2-hop RE & 1-hop RE & 2-hop RE & 1-hop RE \\
\midrule
PrivPetal (official) & $0.090$ & $0.020$ & $0.058$ & $0.015$ \\
\chainnpm (ours) & $0.196$ & $0.105$ & $0.198$ & $0.102$ \\
\bottomrule
\end{tabular}
\end{table}

%=======================================================================
\section{Ablations and efficiency}
\label{app:ablations}

\paragraph{Budget allocation across marginals.}
\label{sec:abl-alloc}
We ablate the equal $\gamma^2$-split by re-weighting the seven
cross-hop marginals (four RI, three RHI) by a factor $w$ relative to the
rest, in \emph{both} directions, $w\in\{0.25,0.5,1,2,4,8\}$. The equal split
is the measured minimum of the resulting curve: at $\varepsilon{=}0.8$ the
median two-hop RE runs $0.93 \to 0.84 \to \mathbf{0.79} \to 1.01 \to 1.19 \to
1.79$ and at $\varepsilon{=}3.2$ $0.10 \to 0.14 \to \mathbf{0.04} \to 0.14
\to 0.21 \to 0.42$ as $w$ increases through that order, so both directions
lose. Concentrating budget on the cross-hop objects starves the
single/adjacent-hop marginals, and join-count utility is set by the noisiest
factor in the join; starving the cross-hop object instead inflates its own
error and erodes $\Delta$-recovery (at $\varepsilon{=}0.8$, $0.44$ at
$w{=}0.25$ and $0.68$ at
$w{=}0.5$, against $0.87$ at $w{=}1$, while for $w\ge1$ recovery stays within
the run-to-run spread, $0.80$--$1.05$). Equal splitting is thus not a
convenience but the empirical optimum of the sweep; workload-driven minimax
allocation is left as future work.
The technical appendix reports this ablation in full, as a
white-box check that the default we adopt is a property of the join rather
than of our tuning.

\paragraph{Other ablations.}
The route-A vs route-B comparison agreed as the design ablation is already
realized by our two synthesis modes: \chainnpm-direct (explicit cross-hop
marginals, no latent mixture $\approx$ route A) vs \chainnpm-mixture
(latent mixture, route B). On the ten-run grid the two modes are indistinguishable at
$\varepsilon{=}3.2$ (Household$\Join$Individual RE $0.060$ direct vs $0.063$
mixture; $\Delta$-recovery $0.98$ vs $0.94$), so on this benchmark the EM
step is not what carries recovery.
\paragraph{Mechanism ablations.}
\emph{(i) Removing the cross-hop marginals.} Replacing the cross-hop
marginals with adjacent one-hop marginals---so that the child attribute is
still measured---collapses $\Delta$-recovery from $0.98$--$1.01$ to
$0.06$--$0.10$ at every budget in direct mode, while 2-hop RE degrades to
the per-FK level ($0.30$--$0.31$, against $0.19$--$0.31$ for the
baselines): the cross-hop marginals are the sole carrier of the multi-hop
signal, and the residual ($\le0.10$) is exactly the observed mediation
URBAN\penalty0$\to$\penalty0OWN\penalty0$\to$\penalty0EMP that per-FK models
also enjoy.
\emph{(ii) EM iterations under a fixed budget.} Sweeping the EM
iteration cap $T\in\{1,2,5,10,20\}$ at fixed $(\varepsilon,\delta)$
leaves \chainnpm-mixture bit-identical---the mixture reaches its fixed
point in one iteration on this benchmark---whereas the PrivLava-style
accounting (fresh budget per iteration, per-iteration noise inflated by
$\sqrt{T}$) degrades Household$\Join$Individual RE from
$0.146$ to $0.313$ at $\varepsilon{=}0.8$ as $T$ grows $1\to20$: budget
dilution, not optimization, governs iterative EM.
\emph{(iii) Types and truncation.}\label{sec:abl-types} $\Delta$-recovery and within-group
RE are nearly invariant to $k\in\{2,3,4,6\}$ and $\tau\in\{4,6,8,12\}$ (at
$\tau{=}4$ the 1-hop RE moves $0.030\to0.031$ and the 2-hop RE by about $5\%$); on this planted
structure the latent types are symmetric and hence unidentifiable, so
recovery is carried entirely by the cross-hop marginals---consistent
with (i).

\subsection{Efficiency}
\label{app:efficiency}
\chainnpm's single budget pass of $Q$ low-order marginals costs
$\mathcal{O}(Q \cdot \sum (n\cdot d))$ time and negligible memory, and
the post-processing stage vectorizes over rows. On TPC-H SF=1, the full
pipeline ($42$ marginals over $536$K orders and $2.1$M lineitems)
completes in $\approx 2.2$ minutes on one CPU; the deep
three-table variant's synthesis takes $\approx 22$--$24$ seconds. PrivPetal's
official implementation needs $1.8$--$2.1$ hours per budget on two A800 GPUs
and $1.7$--$2.3$ hours on one A800 for the deep variant ($\approx 50\times$
wall-clock gap on the standard rail, $\approx 300\times$ on the deep one), and
PrivLava's iterative EM carries the same budget cost.
The gap follows directly from EM being post-processing rather than a
budget-consuming loop.

\paragraph{Cost--effectiveness.} Read as a cost--effect frontier, the same
numbers place \chainnpm{} alone in the region that preserves the effect
(the technical appendix plots the frontier). On the four databases its
retention error leaves the no-cross-table level at both budgets ($0.011$ at
$\varepsilon{=}0.8$ and $0.009$ at $\varepsilon{=}3.2$, against $\approx1.0$
for all five others) while it finishes in under a minute per budget, and
\privpetal{} is dominated at both: $75$--$119\times$ the wall clock for the
same effect level ($1.0006$, i.e.\ no effect retained). On the controlled
benchmark, where all six methods run on one machine, \chainnpm{} is also the
cheapest point in the panel ($2.0$--$2.2$\,s of driver wall clock per run, load included,
against the per-table baseline's $\approx2.5$\,s) and the only series above $\Delta$-recovery
$0.72$. The mechanism thus does not buy effect preservation with wall clock;
the axis on which it pays is query error at low budget
(Section~\ref{P-sec:gap-evidence}).

%=======================================================================

\section{Budget-allocation ablation}
\label{app:alloc}

The submission adopts the equal $\gamma^2$-split of the cross-hop marginals
as its default and reports the sweep in prose; the figure below is the full
ablation of that choice.

\begin{figure*}[t]
\centering
\includegraphics[width=\textwidth]{figs/fig_alloc_sweep3.pdf}
\caption{Budget-allocation ablation (\chainnpm-direct; three runs per point,
segments span the range; two budgets only, $\varepsilon\in\{0.8,3.2\}$, so
the curves carry no six-budget claim). The cross-hop weight $w$ runs over
$\{0.25,0.5,1,2,4,8\}$ on a logarithmic axis. \emph{(a)} Median two-hop query
RE, lowest at $w{=}1$, the equal split, at both budgets. \emph{(b)}
$\Delta$-recovery: unchanged within the run-to-run spread for $w\ge1$
($0.80$--$1.05$) and lower below it. \emph{(c)} The two adjacent-hop classes
(solid region$\times$household, dashed household$\times$individual), both
worse at the two ends of the sweep than at $w{=}1$ and not monotone in
between.}
\label{fig:alloc}
\end{figure*}

%=======================================================================
\section{Ancestor-scale sweep}
\label{app:scale}

The submission quotes the endpoints of this sweep in the conclusion of
\psec{gap-evidence}; the table below reports it in full.

\begin{table*}[t]
\centering
\caption{Ancestor-scale sweep of the controlled benchmark (median over $3$
runs; the same query workload and one stack per column): median two-hop query
relative error by ancestor-table size, for the two \chainnpm{} variants and
the per-table lower bound (\privmrf{}). The smallest scale matches the deep
chain's ancestor
cardinality; there the per-individual marginals saturate (one-hop error $1.0$)
and no method recovers the planted effect, so that column is reported for
completeness rather than compared.}
\label{tab:scale}
\small
\begin{tabular}{@{}lcccc@{}}
\toprule
 & $77$ & $200$ & $700$ & $2500$ \\
\midrule
\multicolumn{5}{@{}l}{$\varepsilon{=}0.8$}\\
\privmrf{} (per-table) & $0.102$ & $0.229$ & $0.241$ & $0.240$ \\
\chainnpm-direct & $1.644$ & $0.791$ & $0.444$ & $0.293$ \\
\chainnpm-mixture & $1.525$ & $0.772$ & $0.438$ & $0.319$ \\
\midrule
\multicolumn{5}{@{}l}{$\varepsilon{=}3.2$}\\
\privmrf{} (per-table) & $0.186$ & $0.209$ & $0.252$ & $0.272$ \\
\chainnpm-direct & $0.426$ & $0.037$ & $0.032$ & $0.015$ \\
\chainnpm-mixture & $0.447$ & $0.114$ & $0.008$ & $0.018$ \\
\midrule
\multicolumn{5}{@{}l}{$\varepsilon{=}8.0$}\\
\privmrf{} (per-table) & $0.212$ & $0.235$ & $0.269$ & $0.294$ \\
\chainnpm-direct & $0.150$ & $0.062$ & $0.043$ & $0.030$ \\
\chainnpm-mixture & $0.219$ & $0.024$ & $0.039$ & $0.030$ \\
\bottomrule
\end{tabular}
\end{table*}

%=======================================================================
\section{Cost--effect frontier}
\label{app:pareto}

The submission states the cost--effectiveness conclusion in prose and
refers here for the frontier itself; the efficiency discussion behind it
is Section~\ref{app:ablations}.

\begin{figure*}[t]
\centering
\includegraphics[width=\textwidth]{figs/fig_pareto.pdf}
\caption{Cost against effect preservation. \emph{(a)} The controlled
benchmark, the same $18$ cells as \pfig{controlled}: $y$ is
$\Delta_{\mathrm{syn}}/\Delta_{\mathrm{real}}$ (median over the three runs;
dotted rule $y{=}0$, dashed rule $y{=}1$) and $x$ the per-cell wall clock the
driver recorded, so all cells come from one machine and one runner. \emph{(b)}
The four real databases: $x$ is the median
over the databases of the per-database median of the synthesis-phase wall
clock (the evaluation is excluded), and $y$ the retention error on the common
basis---\emph{plant\_product} at $t{=}0.316$ and \emph{plant\_xor} at the XOR
target value $0.8$, ten seeds each, pooled per method, dataset and
budget---with the dashed rule at the no-cross-table level.
\privpetal{}'s $x$ is that method's own wall clock on A800 GPUs while the
other five are single-CPU numbers, the asymmetry of
\ref{tab:perbasis-cost}, which is conservative for \chainnpm{}. Each cell's retention error
is relative to that cell's own measured real effect, as everywhere in this
paper; the six methods do not share one realized planted database (the four
per-table and per-FK methods redraw it per seed), and pooling every cell onto
a single denominator moves each mark by at most $0.018$ and leaves
\chainnpm{} the lowest at both budgets. Marker fill encodes the budget
(hollow $\varepsilon{=}0.8$, filled $\varepsilon{=}3.2$, and, in panel (a)
only, half-filled $\varepsilon{=}8.0$), each panel carries its own legend, and
panel (b) uses the palette of
\pfig{four-ds-effect} and~\ref{fig:query-class}.}
\label{fig:pareto}
\end{figure*}

%=======================================================================
\section{Wall-clock cost per basis}
\label{app:timings}

The row-level quota calibration used for every four-database \chainnpm{}
number in the paper is the released flag set
\texttt{--absorb 'pubn,\allowbreak rowleaf,\allowbreak szthresh,\allowbreak
exactmarg,\allowbreak repair,\allowbreak growcap,\allowbreak szquota'};
the default flags reproduce the plain release exactly (verified
field-by-field).

The submission's efficiency section states the cost--accuracy
conclusions and plots the frontier; the per-basis table below is the
timing reference behind its claims. Timings compare within a row, never across
rows, because each row is timed on its own reference chain
implementation.

\begin{table*}[t]
\centering
\caption{Wall-clock cost per budget, stated per basis rather than netted
out---the timing reference behind the submission's efficiency section, with the
chain structure of \ptab{datasets}; each row is timed on its
own reference chain implementation, so timings compare within a row, not
across rows. The \chainnpm{} column
times the end-to-end pipeline on one CPU core (the TPC-H and ACS rows include
restore and CSV output; the four-database rows isolate synthesis; evaluation is
excluded); the
\privpetal{} column is that method's own wall-clock on A800 GPU(s).
``---'': not run on this benchmark.}
\label{tab:perbasis-cost}
\begin{tabular}{@{}llllll@{}}
\toprule
Dataset & Private chain & Scale (rows) & $\varepsilon$ grid & \chainnpm & \privpetal \\
\midrule
\multicolumn{6}{@{}l}{\emph{Established benchmarks}:}\\
Controlled (planted) & R$\leftarrow$H$\leftarrow$I
  & 200 / $\sim$800 / $\sim$2400 & 0.8, 3.2, 8.0 & $\approx$1.1\,s & --- \\
TPC-H SF=1 & Orders, Lineitem
  & 536K / 2.1M & 0.2, 3.2 & 134--136\,s & 1.8--2.1\,h (2\penalty0$\times$\penalty0A800) \\
TPC-H deep & Customer, Orders, Lineitem
  & 150K / 536K / 2.1M & 0.2, 3.2 & $\approx$22--24\,s & 1.7--2.3\,h (1\penalty0$\times$\penalty0A800) \\
ACS PUMS CA'22 & Household, Person
  & 151K / 383K & 0.8, 3.2 & $\approx$20\,s & hours (A800) \\
\midrule
\multicolumn{6}{@{}l}{\emph{Four real databases}:}\\
Financial (PKDD'99) & District$\leftarrow$Account$\leftarrow$Trans
  & 77 / 4.5K / 1.06M & 0.1--3.2 & 6--9\,s & 20--36\,min \\
IMDb & Names$\leftarrow$Principals$\rightarrow$Titles
  & 640K / 1.22M / 110K & 0.1--3.2 & 115--117\,s & 56--106\,min \\
Instacart & Products$\leftarrow$Order-Products$\rightarrow$Orders
  & 38K / 1.52M / 150K & 0.1--3.2 & 141--143\,s & 1.7--3.7\,h \\
MovieLens-1M & Users$\leftarrow$Ratings$\rightarrow$Movies
  & 6.0K / 1.00M / 3.9K & 0.1--3.2 & 46--47\,s & 14--67\,min \\
\bottomrule
\end{tabular}
\end{table*}

\section{Planted-effect results on the four real databases}
\label{app:four-ds}

\textbf{Full planted-effect table.} Table~\ref{tab:four-ds} is the full version of \pfig{four-ds-effect}: retention error per
dataset, protocol, budget and method, with the single-strength note and the
($^{\ddagger}$) marker explained in its caption.

\begin{table*}[t]
\centering
\caption{Planted-effect retention error (median over $10$ seeds; $0$ $=$
effect survives intact, $1$ $=$ lost to the no-cross-table level, $>1$ $=$
distorted) on four real databases under two planting protocols. The deep
chain (financial) reports the contrast $\Delta$; the V-shaped schemas report
the cross-parent cell contrast. Each \emph{plant\_product} cell is the range
over the four planting strengths ($t = 0.1, 0.2, 0.316, 0.4$), so that no
single strength is privileged, except for \privpetal{}, whose planting
protocol varies the privacy budget only and therefore contributes the single
strength
$t{=}0.316$; each \emph{plant\_xor} cell is a single value
at the XOR target value $0.8$, the only strength used.
$^{\ddagger}$the synthetic effect coincides across planting strengths, the
signature of a method whose measured objects do not depend on the planting
(see \pref{sec:four-datasets} for the two benchmark-side caveats that
also produce it).}
\label{tab:four-ds}
\small
\begin{tabular}{@{}llcccccc@{}}
\toprule
Dataset & $\varepsilon$ & \chainnpm & \privmrf & PrivBayes & PB-PGM & PrivLava & \privpetal \\
\midrule
% BEGIN GENERATED tab:four-ds -- do not hand-edit; regenerate with
% results/tables/make_manuscript_tables.py (single source: results/tables/*.tsv)
\multicolumn{8}{@{}l}{\emph{plant\_product}: }\\
financial & 0.8 & 0.415--0.886 & 1.000--1.002 & 0.996--0.999 & 1.002--1.006 & 1.000--1.001$^{\ddagger}$ & 0.998 \\
financial & 3.2 & 0.175--0.346 & 1.000--1.002 & 0.997--0.999 & 1.001--1.004 & 0.996--0.999$^{\ddagger}$ & 1.000 \\
imdb & 0.8 & 0.007--0.011 & 1.003$^{\ddagger}$ & 0.999$^{\ddagger}$ & 1.001$^{\ddagger}$ & 1.011$^{\ddagger}$ & 0.223 \\
imdb & 3.2 & 0.005--0.007 & 0.997$^{\ddagger}$ & 0.998$^{\ddagger}$ & 0.999$^{\ddagger}$ & 1.005$^{\ddagger}$ & 0.055 \\
instacart & 0.8 & 0.007--0.018 & 0.993--0.998$^{\ddagger}$ & 0.997--0.999 & 0.988--0.997$^{\ddagger}$ & 0.996$^{\ddagger}$ & 1.003 \\
instacart & 3.2 & 0.005--0.020 & 1.003--1.013$^{\ddagger}$ & 0.997--0.999 & 1.000--1.001$^{\ddagger}$ & 0.997$^{\ddagger}$ & 1.000 \\
movielens & 0.8 & 0.007--0.021 & 1.003--1.007$^{\ddagger}$ & 1.003--1.007$^{\ddagger}$ & 0.998--0.999$^{\ddagger}$ & 1.002$^{\ddagger}$ & 0.994 \\
movielens & 3.2 & 0.004--0.021 & 0.995--0.998$^{\ddagger}$ & 1.000--1.001$^{\ddagger}$ & 0.999--1.000$^{\ddagger}$ & 0.998$^{\ddagger}$ & 1.010 \\
\midrule
\multicolumn{8}{@{}l}{\emph{plant\_xor}: }\\
financial & 0.8 & 1.614 & 0.998 & 0.992 & 1.012 & 1.002 & 1.049 \\
financial & 3.2 & 0.846 & 0.997 & 0.989 & 1.006 & 0.993 & 1.048 \\
imdb & 0.8 & 0.010 & 1.003 & 0.999 & 1.001 & 1.010 & 0.249 \\
imdb & 3.2 & 0.017 & 0.997 & 0.998 & 0.999 & 1.006 & 0.302 \\
instacart & 0.8 & 0.001 & 1.001 & 1.000 & 1.001 & 1.001 & 1.002 \\
instacart & 3.2 & 0.001 & 0.999 & 1.001 & 1.000 & 1.001 & 1.000 \\
movielens & 0.8 & 0.003 & 1.001 & 1.001 & 1.000 & 1.001 & 1.000 \\
movielens & 3.2 & 0.002 & 0.999 & 1.000 & 1.000 & 0.999 & 0.999 \\
% END GENERATED tab:four-ds
\bottomrule
\end{tabular}
\end{table*}

%=======================================================================
\section{Natural-mode tables and the query-class stratification}
\label{app:natural}

\textbf{Natural-mode query error.} Table~\ref{tab:re-appendix} tabulates the
natural-mode $2$-hop query error of the same four databases, and
Figure~\ref{fig:query-class} gives the same view stratified by query class,
for all six methods and all six budgets. The \chainnpm{} rows are produced by
the row-level quota-calibration arm of \psec{setup}: free post-processing
that spends no budget and changes no released marginal. With it \chainnpm{} is
the most accurate method on the $2$-hop class on $22$ of the $24$
dataset--budget cells (financial ties \privmrf{} at
$\varepsilon{=}0.4$ and $0.8$, the two gaps $\le 0.02$)---the deep
chain's error collapses from $4.34$ at $\varepsilon{=}0.1$ under the released
arm to $0.28$, and the V-shaped schemas lead their per-table baselines
throughout (Table~\ref{tab:re-appendix}); on the class-stratified medians of
Figure~\ref{fig:query-class} it leads every cross-table class ($2$-hop,
both $1$-hop spans, and the two-parent class) at every budget, and the
zero-hop class it was weakest in no longer shows the released arm's
order-of-magnitude gap on the V-shaped schemas. Natural-mode query accuracy
is therefore not a liability we qualify away: the correlation the per-FK
family cannot represent is preserved (Table~\ref{tab:four-ds}), and
carrying
it no longer costs query accuracy, as the submission's planted-effect
figure (\pfig{four-ds-effect}) and Table~\ref{tab:re-appendix}
show.

\begin{table*}[t]
\centering
\caption{Natural mode (no planting): $2$-hop large-query median relative
error over $10$ seeds, i.e.\ the cross-hop class
(\textsc{District}$\times$\textsc{Trans} for financial; the two-parent class
for the V-shaped schemas). Lower is better. \chainnpm{} rows use the
row-level quota-calibration arm of \psec{setup} (the released
default is bit-identical to the plain release); the arm spends no budget and
changes no marginal, and it removes the released version's natural-mode
query-accuracy gap on all four databases. The table is reported in full
because the ranking is now part of the result, not a caveat to it.}
\label{tab:re-appendix}
\small
\begin{tabular}{@{}llcccccc@{}}
\toprule
Method & Dataset & $\varepsilon{=}0.1$ & $\varepsilon{=}0.2$ & $\varepsilon{=}0.4$ & $\varepsilon{=}0.8$ & $\varepsilon{=}1.6$ & $\varepsilon{=}3.2$ \\
\midrule
% BEGIN GENERATED tab:re-appendix -- do not hand-edit; regenerate with
% results/tables/make_manuscript_tables.py (single source: results/tables/*.tsv)
\chainnpm & financial & 0.2774 & 0.2123 & 0.1868 & 0.0977 & 0.0969 & 0.0665 \\
PrivMRF & financial & 0.2875 & 0.2183 & 0.1698 & 0.0969 & 0.1357 & 0.1555 \\
PrivBayes & financial & 0.5180 & 0.4960 & 0.4017 & 0.1383 & 0.1734 & 0.1186 \\
PB-PGM & financial & 1.0000 & 0.8510 & 0.6907 & 0.2802 & 0.2667 & 0.1367 \\
PrivLava & financial & 0.3402 & 0.3798 & 0.3429 & 0.3686 & 0.3304 & 0.3430 \\
\privpetal & financial & 0.2921 & 0.5083 & 0.7317 & 0.8421 & 0.9166 & 0.9430 \\
\midrule
\chainnpm & imdb & 0.1323 & 0.1286 & 0.1172 & 0.1157 & 0.1133 & 0.1106 \\
PrivMRF & imdb & 0.2050 & 0.2029 & 0.2049 & 0.2004 & 0.1954 & 0.1948 \\
PrivBayes & imdb & 0.2205 & 0.2221 & 0.2114 & 0.2028 & 0.1983 & 0.1974 \\
PB-PGM & imdb & 0.2230 & 0.2132 & 0.2101 & 0.2157 & 0.2024 & 0.2072 \\
PrivLava & imdb & 0.9992 & 0.9992 & 0.9993 & 0.9992 & 0.9991 & 0.9993 \\
\privpetal & imdb & 0.6410 & 0.3064 & 0.3127 & 0.2925 & 0.3318 & 0.3642 \\
\midrule
\chainnpm & instacart & 0.0317 & 0.0275 & 0.0281 & 0.0254 & 0.0260 & 0.0254 \\
PrivMRF & instacart & 0.4912 & 0.4984 & 0.4780 & 0.5037 & 0.4915 & 0.4791 \\
PrivBayes & instacart & 0.6624 & 0.5591 & 0.5011 & 0.4919 & 0.4662 & 0.4709 \\
PB-PGM & instacart & 0.5016 & 0.4854 & 0.4692 & 0.4766 & 0.4821 & 0.4768 \\
PrivLava & instacart & 0.9044 & 0.9008 & 0.9083 & 0.9198 & 0.9337 & 0.9534 \\
\privpetal & instacart & 1.2073 & 1.0324 & 0.4536 & 0.4443 & 0.4428 & 0.4383 \\
\midrule
\chainnpm & movielens & 0.2127 & 0.1690 & 0.1590 & 0.1345 & 0.1291 & 0.1301 \\
PrivMRF & movielens & 0.3894 & 0.3589 & 0.3521 & 0.3444 & 0.3409 & 0.3314 \\
PrivBayes & movielens & 0.4087 & 0.3736 & 0.3620 & 0.3515 & 0.3480 & 0.3322 \\
PB-PGM & movielens & 0.4679 & 0.3886 & 0.3829 & 0.3493 & 0.3503 & 0.3353 \\
PrivLava & movielens & 0.9805 & 0.9827 & 0.9819 & 0.9686 & 0.9742 & 0.9634 \\
\privpetal & movielens & 0.6667 & 0.7739 & 0.8181 & 0.8007 & 0.8106 & 0.8166 \\
% END GENERATED tab:re-appendix
\bottomrule
\end{tabular}
\end{table*}

\begin{figure*}[t]
\centering
\includegraphics[width=\textwidth]{figs/fig_query_class_sweep.pdf}
\caption{The natural-mode runs stratified by query class: $2$-hop
conjunctive, $1$-hop predicate, $0$-hop marginal, and size-distribution
total variation. A class value averages that schema's member metrics
arithmetically (the chain's two adjacent-hop objects, a V-shaped schema's
two or three one-way marginals), then the $10$ seeds are reduced to their
median, then the four
databases to theirs. Each curve is therefore the median over the datasets on
which \emph{all six} methods report that class, so no curve inside a panel is
aggregated over a different basis than its neighbours; the size-distribution
panel consequently reduces to the deep chain, the only database where all six
report it. Vertical scales differ between panels and must not be compared
across them.}
\label{fig:query-class}
\end{figure*}

%=======================================================================
\section{Positioning of representative mechanisms}
\label{app:position}

The paper's related-work discussion carries the positioning; the table
below is its tabular form.

\begin{table*}[t]
\centering
\caption{Positioning of representative DP relational synthesis
mechanisms. \emph{Cross-table object}: the statistical object each
mechanism explicitly models across foreign keys; \emph{budget pattern}:
iterative budget consumption vs.\ a single measurement pass followed by
free post-processing; \emph{evaluated private depth}: the deepest
private join path measured in the original evaluation; ``---'': not
reported in the cited work.}
\label{tab:position}
\small
\begin{tabular}{@{}p{2.1cm}p{4.7cm}p{3.3cm}p{2.7cm}p{2.7cm}@{}}
\toprule
Method & Cross-table object & Budget pattern & Evaluated private depth & Runtime order \\
\midrule
\privlava (SIGMOD'23) & per-FK latent models (EM); cross-hop only via
estimated $\hat Z$ ($\le 16$ values) & iterative: each EM iteration
consumes fresh budget & 1 private FK (upstream tables public) &
infeasible at EP scale \\
\privpetal (SIGMOD'25) & per-FK permutation marginals; stages condition
on synthetic parents; no joint object across two FKs & per-stage split
(e.g., $0.5/0.25/0.25$) & 1 FK hop (query classes) & 1.7--2.3\,h per
budget (A800) \\
Adapted single-table \cite{alimohammadi2024adapting} & pairwise
cross-table marginals only & --- & --- & --- \\
\chainnpm (ours) & joint chain marginals incl.\ cross-hop
$R_0{\times}R_2$; latent mixture refines only & single pass; EM as free
post-processing & 3 private tables & $\approx$22--24\,s (1 CPU) \\
\bottomrule
\end{tabular}
\end{table*}

\section{The three real-benchmark rails}
\label{app:rails}

Figure~\ref{fig:legacy-rails} collects the three real-benchmark rails of
Sections~\ref{app:tpch-rails}--\ref{app:acs-rail}---standard TPC-H, the deep
three-table variant, and ACS PUMS---in one place: each panel plots query
relative error against the budgets that rail was actually run at, with one
line per method and per hop class. The rails were run at two budgets each, so
this summary is deliberately sparser than the four-database sweep of
Table~\ref{tab:re-appendix} and the class-stratified view below, and is meant to
be read
together with the submission's tables for the same rails, not instead of
them.

\begin{figure*}[t]
\centering
\includegraphics[width=\textwidth]{figs/fig_legacy_rails.pdf}
\caption{The three real-benchmark rails on one page: standard TPC-H (two
private tables), the deep three-table variant, and ACS PUMS. Each panel uses
that rail's own budgets ($\varepsilon\in\{0.2,3.2\}$ for the two TPC-H
variants, $\{0.8,3.2\}$ for ACS) and plots query relative error; solid lines
are the deepest hop that rail reports (two-hop for TPC-H and ACS, three-hop for
the deep variant), and the dashed line is ACS's one-hop class. Values are the
median relative errors tabulated in \ref{tab:tpch} and Table~\ref{tab:acs}.}
\label{fig:legacy-rails}
\end{figure*}

% The moved prose cites the mechanisms it compares against, so the companion
% needs the same bibliography as the paper (BibTeX lists only what is cited).
\bibliographystyle{ACM-Reference-Format}
\bibliography{refs}

\end{document}